\chapter{2015年广东卷（文）}

\section{单选题}

\begin{problem}
若集合 \(M = \{-1, 1\}\)，\(N = \{-2, 1, 0\}\)，则 \(M \cap N =\)（\quad）

\choices
    {\(\{0, -1\}\)}
    {\(\{1\}\)}
    {\(\{0\}\)}
    {\(\{-1, 1\}\)}
\end{problem}

\begin{answer}
B
\end{answer}

\begin{solution}
交集中的元素必须同时属于 \(M\) 和 \(N\). 集合 \(M\) 的元素为 \(-1\)，\(1\)，其中只有 \(1\) 也属于 \(N\)，所以 \(M\cap N=\{1\}\)，故选 B.
\end{solution}

\begin{problem}
已知 \(\i\) 是虚数单位，则复数 \((1 + \i)^{2}\) =（\quad）

\choices
    {\(2\i\)}
    {\(-2\i\)}
    {\(2\)}
    {\(-2\)}
\end{problem}

\begin{answer}
A
\end{answer}

\begin{solution}
由 \(\i^2=-1\)，直接展开得 \((1+\i)^2=1+2\i+\i^2=1+2\i-1=2\i\)，故选 A.
\end{solution}

\begin{problem}
下列函数中，既不是奇函数，也不是偶函数的是（\quad）

\choices
    {\(y = x + \sin 2x\)}
    {\(y = x^{2} - \cos x\)}
    {\(y = 2^{x} + \frac{1}{2^{x}}\)}
    {\(y = x^{2} + \sin x\)}
\end{problem}

\begin{answer}
D
\end{answer}

\begin{solution}
四个函数的定义域都是 \(\R\). 对于 A，\(f(-x)=-x-\sin 2x=-f(x)\)，所以它是奇函数；对于 B，\(f(-x)=x^2-\cos x=f(x)\)，所以它是偶函数. 对于 C，\(f(x)=2^x+2^{-x}\)，从而 \(f(-x)=2^{-x}+2^x=f(x)\)，所以它是偶函数. 对于 D，\(f(-x)=x^2-\sin x\)，既不恒等于 \(f(x)\)，也不恒等于 \(-f(x)\)，因此 D 中的函数既不是奇函数，也不是偶函数，故选 D.
\end{solution}

\begin{problem}
若变量 \(x\)，\(y\) 满足约束条件 \(\begin{cases}
x + 2y \leqslant 2,\\
x + y \geqslant 0,\\
x \leqslant 4,
\end{cases}\) 则 \(z = 2x + 3y\) 的最大值为（\quad）

\choices
    {\(2\)}
    {\(5\)}
    {\(8\)}
    {\(10\)}
\end{problem}

\begin{answer}
B
\end{answer}

\begin{solution}
三条边界直线的交点分别为：\(x+2y=2\) 与 \(x+y=0\) 的交点 \((-2,2)\)，\(x+2y=2\) 与 \(x=4\) 的交点 \((4,-1)\)，以及 \(x+y=0\) 与 \(x=4\) 的交点 \((4,-4)\). 约束条件表示的可行域是这三个点构成的三角形，线性目标函数的最大值在顶点处取得. 于是
\(z\big|_{(-2,2)}=2\)，\(z\big|_{(4,-1)}=5\)，\(z\big|_{(4,-4)}=-4\).
因此 \(z\) 的最大值为 \(5\)，故选 B.
\end{solution}

\begin{problem}
设 \(\triangle ABC\) 的内角 \(A\)，\(B\)，\(C\) 的对边分别为 \(a\)，\(b\)，\(c\). 若 \(a = 2\)，\(c = 2\sqrt{3}\)，\(\cos A = \frac{\sqrt{3}}{2}\) 且 \(b < c\)，则 \(b =\)（\quad）

\choices
    {\(3\)}
    {\(2\sqrt{2}\)}
    {\(2\)}
    {\(\sqrt{3}\)}
\end{problem}

\begin{answer}
C
\end{answer}

\begin{solution}
由余弦定理，\(a^2=b^2+c^2-2bc\cos A\). 代入题设数据，得
\[
4=b^2+12-2\cdot b\cdot2\sqrt{3}\cdot\frac{\sqrt{3}}{2}=b^2+12-6b
\]
即 \(b^2-6b+8=0\)，所以 \(b=2\) 或 \(b=4\). 由于 \(b<c=2\sqrt{3}<4\)，只能取 \(b=2\)，故选 C.
\end{solution}

\begin{problem}
若直线 \(l_{1}\) 和 \(l_{2}\) 是异面直线，\(l_{1}\) 在平面 \(\alpha\) 内，\(l_{2}\) 在平面 \(\beta\) 内，\(l\) 是平面 \(\alpha\) 与平面 \(\beta\) 的交线. 则下列命题正确的是（\quad）

\choices
    {\(l\) 与 \(l_{1}\)、\(l_{2}\) 都不相交}
    {\(l\) 与 \(l_{1}\)、\(l_{2}\) 都相交}
    {\(l\) 至多与 \(l_{1}\)、\(l_{2}\) 中的一条相交}
    {\(l\) 至少与 \(l_{1}\)、\(l_{2}\) 中的一条相交}
\end{problem}

\begin{answer}
D
\end{answer}

\begin{solution}
反设 \(l\) 与 \(l_1\)、\(l_2\) 都不相交. 因为 \(l_1\) 与 \(l\) 同在平面 \(\alpha\) 内，所以 \(l_1\parallel l\)；同理，\(l_2\parallel l\). 于是 \(l_1\parallel l_2\)，这与 \(l_1\) 和 \(l_2\) 是异面直线矛盾. 因此 \(l\) 至少与 \(l_1\)、\(l_2\) 中的一条相交，故选 D.
\end{solution}

\begin{problem}
已知 5 件产品中有 2 件次品，其余为合格品，现从 5 件产品中任取 2 件，恰有一件次品的概率为（\quad）

\choices
    {\(0.4\)}
    {\(0.6\)}
    {\(0.8\)}
    {\(1\)}
\end{problem}

\begin{answer}
B
\end{answer}

\begin{solution}
从 5 件产品中任取 2 件的基本事件数为 \(\mathrm C_5^2=10\). 从 2 件次品中取 1 件、从 3 件合格品中取 1 件的有利事件数为 \(\mathrm C_2^1\mathrm C_3^1=6\). 因此所求概率为
\[
P=\frac{\mathrm C_2^1\mathrm C_3^1}{\mathrm C_5^2}=\frac{6}{10}=0.6
\]
故选 B.
\end{solution}

\begin{problem}
已知椭圆 \(\frac{x^2}{25} + \frac{y^2}{m^2} = 1\)（\(m > 0\)）的左焦点为 \(F_{1}(-4,0)\)，则 \(m =\)（\quad）

\choices
    {\(2\)}
    {\(3\)}
    {\(4\)}
    {\(9\)}
\end{problem}

\begin{answer}
B
\end{answer}

\begin{solution}
椭圆的长半轴为 \(a=5\)，焦距的一半为 \(c=4\). 由椭圆关系 \(a^2=b^2+c^2\)，且短半轴 \(b=m>0\)，得 \(m^2=25-16=9\)，所以 \(m=3\)，故选 B.
\end{solution}

\begin{problem}
在平面直角坐标系 \(xOy\) 中，已知四边形 \(ABCD\) 是平行四边形，\(\overrightarrow{AB} = \left(1, -2\right)\)，\(\overrightarrow{AD} = \left(2, 1\right)\)，则 \(\overrightarrow{AD} \cdot \overrightarrow{AC} =\)（\quad）

\choices
    {\(5\)}
    {\(4\)}
    {\(3\)}
    {\(2\)}
\end{problem}

\begin{answer}
A
\end{answer}

\begin{solution}
平行四边形的对角线满足 \(\overrightarrow{AC}=\overrightarrow{AB}+\overrightarrow{AD}\)，所以 \(\overrightarrow{AC}=(3,-1)\). 于是
\[
\overrightarrow{AD}\cdot\overrightarrow{AC}=\left(2,1\right)\cdot\left(3,-1\right)=2\times3+1\times(-1)=5
\]
故选 A.
\end{solution}

\begin{problem}
若集合 \(E = \{(p, q, r, s) \mid 0 \leqslant p < s \leqslant 4, 0 \leqslant q < s \leqslant 4, 0 \leqslant r < s \leqslant 4\text{ 且 }p, q, r, s \in \N\}\)，\(F = \{(t, u, v, w) \mid 0 \leqslant t < u \leqslant 4, 0 \leqslant v < w \leqslant 4\text{ 且 }t, u, v, w \in \N\}\)，用 \(\operatorname{card}(X)\) 表示集合 \(X\) 中元素个数，则 \(\operatorname{card}(E) + \operatorname{card}(F) =\)（\quad）

\choices
    {\(200\)}
    {\(150\)}
    {\(100\)}
    {\(50\)}
\end{problem}

\begin{answer}
A
\end{answer}

\begin{solution}
固定 \(s\). 当 \(s=1\)，\(2\)，\(3\)，\(4\) 时，\(p\)，\(q\)，\(r\) 分别有 \(1\)，\(2\)，\(3\)，\(4\) 种取法，故
\[
\operatorname{card}(E)=1^3+2^3+3^3+4^3=100
\]
在集合 \(F\) 中，\((t,u)\) 是 \(0\)，\(1\)，\(2\)，\(3\)，\(4\) 中两个不同数按从小到大排列的有序对，因此有 \(\mathrm C_5^2=10\) 种取法；\((v,w)\) 同样有 10 种取法，故 \(\operatorname{card}(F)=10\times10=100\). 所以 \(\operatorname{card}(E)+\operatorname{card}(F)=200\)，故选 A.
\end{solution}

\section{填空题}

\begin{problem}
不等式 \(-x^{2} - 3x + 4 > 0\) 的解集为\fillinblank{}.（用区间表示）
\end{problem}

\begin{answer}
\(\left(-4,1\right)\)
\end{answer}

\begin{solution}
不等式等价于 \(x^2+3x-4<0\)，而 \(x^2+3x-4=(x+4)(x-1)\). 抛物线开口向上，所以乘积小于零时 \(-4<x<1\)，解集为 \(\left(-4,1\right)\).
\end{solution}

\begin{problem}
已知样本数据 \(x_{1}\)，\(x_{2}\)，\(\cdots\)，\(x_{n}\) 的均值 \(\overline{x} = 5\)，则样本数据 \(2x_{1} + 1\)，\(2x_{2} + 1\)，\(\cdots\)，\(2x_{n} + 1\) 的均值为\fillinblank{}.
\end{problem}

\begin{answer}
11
\end{answer}

\begin{solution}
设新样本的均值为 \(\overline{y}\)，则
\[
\overline{y}=\frac{1}{n}\sum_{i=1}^{n}(2x_i+1)=2\left(\frac{1}{n}\sum_{i=1}^{n}x_i\right)+1=2\overline{x}+1=2\times5+1=11
\]
所以所填数为 \(11\).
\end{solution}

\begin{problem}
若三个正数 \(a\)，\(b\)，\(c\) 成等比数列，其中 \(a = 5 + 2\sqrt{6}\)，\(c = 5 - 2\sqrt{6}\)，则 \(b =\)\fillinblank{}.
\end{problem}

\begin{answer}
1
\end{answer}

\begin{solution}
三个正数成等比数列，所以中项满足 \(b^2=ac\). 代入得 \(b^2=(5+2\sqrt6)(5-2\sqrt6)=25-24=1\). 因为 \(b>0\)，所以 \(b=1\).
\end{solution}

\section{选考题：填空题}

\begin{problem}
在平面直角坐标系 \(xOy\) 中，以原点 \(O\) 为极点，\(x\) 轴的正半轴为极轴建立极坐标系，曲线 \(C_1\) 的极坐标方程为 \(\rho (\cos \theta + \sin \theta) = -2\)，曲线 \(C_2\) 的参数方程为 \(\begin{cases}
x = t^2,\\
y = 2\sqrt{2}t,
\end{cases}\)（\(t\) 为参数），则 \(C_1\) 与 \(C_2\) 交点的直角坐标为\fillinblank{}.
\end{problem}

\begin{answer}
\(\left(2,-4\right)\)
\end{answer}

\begin{solution}
由极坐标与直角坐标的关系 \(x=\rho\cos\theta\)，\(y=\rho\sin\theta\)，曲线 \(C_1\) 的方程化为 \(x+y=-2\). 由参数方程消去 \(t\)：\(y^2=8t^2=8x\). 联立得
\(x+y=-2\)，\(y^2=8x\).
将 \(x=-2-y\) 代入第二式，得 \(y^2=8(-2-y)\)，即 \((y+4)^2=0\). 于是 \(y=-4\)，\(x=2\)，交点坐标为 \(\left(2,-4\right)\).
\end{solution}

\begin{problem}
如图，\(AB\) 为圆 \(O\) 的直径，\(E\) 为 \(AB\) 延长线上一点，过 \(E\) 作圆 \(O\) 的切线，切点为 \(C\)，过 \(A\) 作直线 \(EC\) 的垂线，垂足为 \(D\). 若 \(AB = 4\)，\(CE = 2\sqrt{3}\)，则 \(AD =\)\fillinblank{}.
\bitmapfigure[max width=0.46\linewidth]{img/2015/guangdong_liberal/q15_fig1.png}
\end{problem}

\begin{answer}
3
\end{answer}

\begin{solution}
设 \(EB=u>0\)，则 \(EA=u+4\). 由切割线定理，\(EC^2=EB\cdot EA\)，所以 \(12=u(u+4)\)，即 \((u-2)(u+6)=0\). 因 \(u>0\)，得 \(u=2\)，从而 \(EA=6\)，\(EO=EB+BO=2+2=4\).

因为 \(OC\perp EC\)，且 \(AD\perp EC\)，所以 \(AD\parallel OC\)，从而直角三角形 \(EAD\) 与 \(EOC\) 相似. 于是
\[
\frac{AD}{OC}=\frac{EA}{EO}=\frac{6}{4}
\]
而 \(OC=\frac{AB}{2}=2\)，故 \(AD=3\).
\end{solution}

\section{解答题}

\begin{problem}
已知 \(\tan \alpha = 2\).
\begin{enumerate}
\item 求 \( \tan\left(\alpha+\frac{\pi}{4}\right) \) 的值；
\item 求 \( \frac{\sin 2\alpha}{\sin^{2}\alpha + \sin \alpha \cos \alpha - \cos 2\alpha - 1} \) 的值.
\end{enumerate}
\end{problem}

\begin{solution}
\begin{enumerate}
\item 由正切的和角公式，且 \(1-\tan\alpha\neq0\)，得
\[
\tan\left(\alpha+\frac{\pi}{4}\right)=\frac{\tan\alpha+\tan\frac{\pi}{4}}{1-\tan\alpha\tan\frac{\pi}{4}}=\frac{2+1}{1-2}=-3
\]

\item 由 \(\tan\alpha=2\)，可设 \(\sin\alpha=2\cos\alpha\). 结合 \(\sin^2\alpha+\cos^2\alpha=1\)，得 \(5\cos^2\alpha=1\)，即 \(\cos^2\alpha=\frac15\). 因此 \(\sin2\alpha=2\sin\alpha\cos\alpha=4\cos^2\alpha=\frac45\)，且 \(\sin^2\alpha+\sin\alpha\cos\alpha=6\cos^2\alpha=\frac65\)，\(\cos2\alpha=\cos^2\alpha-\sin^2\alpha=-3\cos^2\alpha=-\frac35\). 原分式的分母为 \(\frac65-\left(-\frac35\right)-1=\frac45\)，所以原分式的值为 \(\frac{4/5}{4/5}=1\).
\end{enumerate}
\end{solution}

\begin{problem}
某城市 100 户居民的月平均用电量（单位：\(\text{度}\)），\([160, 180)\)，\([180, 200)\)，\([200, 220)\)，\([220, 240)\)，\([240, 260)\)，\([260, 280)\)，\([280, 300)\) 分组的频率分布直方图如图.

\begin{enumerate}
\item 求直方图中 \(x\) 的值；
\item 求月平均用电量的众数和中位数；
\item 在月平均用电量为 \([220, 240)\)，\([240, 260)\)，\([260, 280)\)，\([280, 300)\) 的四组用户中，用分层抽样的方法抽取 11 户居民，则月平均用电量在 \([220, 240)\) 的用户中应抽取多少户？
\end{enumerate}
\bitmapfigure[max width=0.52\linewidth]{img/2015/guangdong_liberal/q17_fig1.png}
\end{problem}

\begin{solution}
\begin{enumerate}
\item 各组组距均为 \(20\)，频率等于组距乘以频率组距，所以由直方图总面积为 1，得
\[
20\left(0.002+0.0095+0.011+0.0125+x+0.005+0.0025\right)=1
\]
解得 \(x=0.0075\).

\item 直方图中高度最大的组为 \([220,240)\)，所以它是众数组. 用众数组的组中值估计众数，得众数为 \(\frac{220+240}{2}=230\). 各组的频数依次为 \(4\)，\(19\)，\(22\)，\(25\)，\(15\)，\(10\)，\(5\)，累计频数为 \(4\)，\(23\)，\(45\)，\(70\)，\(85\)，\(95\)，\(100\). 第 50 个数据落在 \([220,240)\) 内，因此中位数为
\[
220+20\times\frac{50-45}{25}=224
\]

\item 四组用户的频率之和为 \(0.25+0.15+0.10+0.05=0.55\)，其中 \([220,240)\) 组的频率为 \(0.25\). 按分层抽样比例，应抽取
\[
11\times\frac{0.25}{0.55}=5
\]
户.
\end{enumerate}
\end{solution}

\begin{problem}
如图，三角形 \(PDC\) 所在的平面与长方形 \(ABCD\) 所在的平面垂直，\(PD = PC = 4\)，\(AB = 6\)，\(BC = 3\).
\begin{enumerate}
\item 证明： \(BC \parallel\) 平面 \(PDA\)；
\item 证明： \(BC \perp PD\)；
\item 求点 \(C\) 到平面 \(PDA\) 的距离.
\end{enumerate}
\FigureLayoutDeclare{img/2015/guangdong_liberal/q18_fig1.png}{10.9cm}{15.839bp 62.875bp 103.191bp 1.92bp}
\bitmapfigure[max width=0.40\linewidth]{img/2015/guangdong_liberal/q18_fig1.png}
\end{problem}

\begin{solution}
取长方形所在平面为 \(z=0\)，建立空间直角坐标系，使 \(A=(0,0,0)\)，\(B=(6,0,0)\)，\(D=(0,3,0)\)，\(C=(6,3,0)\). 由于平面 \(PDC\) 与 \(z=0\) 垂直且交线为 \(DC\)，可令 \(P=(3,3,h)\). 这里 \(P\) 的横坐标为 3 是由 \(PD=PC\) 得到的；再由 \(PD=4\) 得 \(9+h^2=16\)，所以 \(h^2=7\).

平面 \(PDA\) 过原点，且 \(\overrightarrow{AD}=(0,3,0)\)，\(\overrightarrow{AP}=(3,3,h)\)，故其一个法向量为 \((3h,0,-9)\)，平面方程为
\[
hx-3z=0
\]

\begin{enumerate}
\item \(\overrightarrow{BC}=(0,3,0)=\overrightarrow{AD}\)，所以 \(BC\parallel AD\). 又点 \(B=(6,0,0)\) 不满足平面方程 \(hx-3z=0\)，故直线 \(BC\) 不在平面 \(PDA\) 内，从而 \(BC\parallel\) 平面 \(PDA\).

\item \(\overrightarrow{DP}=(3,0,h)\)，于是 \(\overrightarrow{BC}\cdot\overrightarrow{DP}=0\times3+3\times0+0\times h=0\). 因此 \(BC\perp PD\).

\item 点 \(C=(6,3,0)\) 到平面 \(hx-3z=0\) 的距离为
\[
\frac{|6h|}{\sqrt{h^2+9}}=\frac{6\sqrt7}{\sqrt{7+9}}=\frac{3\sqrt7}{2}
\]
所以点 \(C\) 到平面 \(PDA\) 的距离为 \(\dfrac{3\sqrt7}{2}\).
\end{enumerate}
\end{solution}

\begin{problem}
设数列 \(\{a_{n}\}\) 的前 \(n\) 项和为 \(S_{n}\)，\(n \in \N^{*}\)，已知 \(a_{1} = 1\)，\(a_{2} = \frac{3}{2}\)，\(a_{3} = \frac{5}{4}\)，且当 \(n \geqslant 2\) 时，\(4S_{n+2} + 5S_{n} = 8S_{n+1} + S_{n-1}\).
\begin{enumerate}
\item 求 \( a_{4} \) 的值.
\item 证明： \( \left\{a_{n+1}-\frac{1}{2}a_{n}\right\} \) 为等比数列；
\item 求数列 \( \{a_{n}\} \) 的通项公式.
\end{enumerate}
\end{problem}

\begin{solution}
\begin{enumerate}
\item 由 \(S_1=1\)，\(S_2=1+\frac32=\frac52\)，\(S_3=1+\frac32+\frac54=\frac{15}{4}\). 在已知递推关系中取 \(n=2\)，得
\[
4S_4+5S_2=8S_3+S_1
\]
又 \(S_4=S_3+a_4\)，所以 \(4\left(\frac{15}{4}+a_4\right)+5\cdot\frac52=8\cdot\frac{15}{4}+1\)，解得 \(a_4=\frac78\).

\item 将 \(S_{n+2}=S_{n+1}+a_{n+2}\)、\(S_{n+1}=S_n+a_{n+1}\)、\(S_n=S_{n-1}+a_n\) 代入已知关系，得
\(4a_{n+2}-4a_{n+1}+a_n=0\)，\(n\geqslant2\).
令 \(b_n=a_{n+1}-\frac12a_n\)，则
\[
b_{n+1}=a_{n+2}-\frac12a_{n+1}=\frac12a_{n+1}-\frac14a_n=\frac12b_n
\]
其中 \(b_1=a_2-\frac12a_1=1\)，\(b_2=a_3-\frac12a_2=\frac12=\frac12b_1\). 上式对 \(n\geqslant2\) 成立，因此 \(\{b_n\}\) 是首项为 1、公比为 \(\frac12\) 的等比数列，即 \(\left\{a_{n+1}-\frac12a_n\right\}\) 为等比数列.

\item 由上一步，\(a_{n+1}-\frac12a_n=\left(\frac12\right)^{n-1}\). 两边乘以 \(2^n\)，并令 \(c_n=2^{n-1}a_n\)，得 \(c_{n+1}-c_n=2\). 又 \(c_1=a_1=1\)，所以 \(c_n=1+2(n-1)=2n-1\). 因此
\[
a_n=\frac{c_n}{2^{n-1}}=(2n-1)\left(\frac12\right)^{n-1}
\]
\end{enumerate}
\end{solution}

\begin{problem}
已知过原点的动直线 \(l\) 与圆 \(C_1: x^2 + y^2 - 6x + 5 = 0\) 相交于不同的两点 \(A\)，\(B\).
\begin{enumerate}
\item 求圆 \(C_1\) 的圆心坐标；
\item 求线段 \(AB\) 的中点 \(M\) 的轨迹 \(C\) 的方程；
\item 是否存在实数 \(k\)，使得直线 \(L: y = k(x - 4)\) 与曲线 \(C\) 只有一个交点？若存在，求出 \(k\) 的取值范围；若不存在，说明理由.
\end{enumerate}
\end{problem}

\begin{solution}
\begin{enumerate}
\item 将圆方程配方，得 \((x-3)^2+y^2=4\)，所以圆心为 \((3,0)\)，半径为 2.

\item 记圆心为 \(I=(3,0)\). 圆心到弦 \(AB\) 的垂线平分弦 \(AB\)，所以 \(IM\perp OM\)，其中 \(O=(0,0)\). 设 \(M=(x,y)\)，则
\[
(x-3,y)\cdot(x,y)=0
\]
即 \(x^2+y^2-3x=0\). 还需限制直线 \(l\) 与圆有两个不同交点. 因为 \(IM\) 是圆心到直线 \(l\) 的距离，所以必须有 \(IM<2\). 在轨迹方程上，\(IM^2=(x-3)^2+y^2=9-3x\)，故 \(9-3x<4\)，即 \(x>\frac53\). 因此
\(C:\ x^2+y^2-3x=0\)，\(\frac53<x\leqslant3\).

\item 曲线 \(C\) 关于 \(x\) 轴对称，先讨论 \(k\geqslant0\). 下半圆弧上的点满足 \(y=-\sqrt{x(3-x)}\)，而该点在直线 \(L\) 上时
\[
k=\frac{\sqrt{x(3-x)}}{4-x}=:g(x)
\]
其中 \(\frac53<x\leqslant3\).
计算 \(g^2(x)=\frac{x(3-x)}{(4-x)^2}\)，得
\[
\left(g^2(x)\right)'=\frac{12-5x}{(4-x)^3}
\]
所以 \(g(x)\) 在 \(\left(\frac53,\frac{12}{5}\right]\) 上递增，在 \(\left[\frac{12}{5},3\right]\) 上递减，并且
\(\lim_{x\to\frac53+}g(x)=\frac{2\sqrt5}{7}\)，\(g\left(\frac{12}{5}\right)=\frac34\)，\(g(3)=0\)
左端点 \(x=\frac53\) 不属于轨迹，因此当 \(0\leqslant k\leqslant\frac{2\sqrt5}{7}\) 时只有递减支上的一个交点；当 \(k=\frac34\) 时在极大值点有一个交点；当 \(\frac{2\sqrt5}{7}<k<\frac34\) 时有两个交点，其他 \(k\geqslant0\) 没有交点. 利用关于 \(x\) 轴的对称性，得到
\[
k\in\left[-\frac{2\sqrt5}{7},\frac{2\sqrt5}{7}\right]\cup\left\{-\frac34,\frac34\right\}
\]
所以这样的实数 \(k\) 存在，且取值范围如上.
\end{enumerate}
\end{solution}

\begin{problem}
设 \(a\) 为实数，函数 \(f(x) = (x - a)^2 + |x - a| - a(a - 1)\).
\begin{enumerate}
\item 若 \( f(0) \leqslant 1 \)，求 \( a \) 的取值范围；
\item 讨论 \( f(x) \) 的单调性；
\item 当 \(a \geqslant 2\) 时，讨论 \(f(x) + \frac{4}{x}\) 在区间 \((0, +\infty)\) 内的零点个数.
\end{enumerate}
\end{problem}

\begin{solution}
\begin{enumerate}
\item \(f(0)=a^2+|a|-a(a-1)=|a|+a\). 当 \(a\leqslant0\) 时，\(|a|+a=0\leqslant1\)；当 \(a>0\) 时，\(|a|+a=2a\leqslant1\)，即 \(a\leqslant\frac12\). 综上，\(a\leqslant\frac12\).

\item 当 \(x<a\) 时，\(|x-a|=a-x\)，所以
\[
f(x)=x^2-(2a+1)x+2a
\]
其导数为 \(2x-(2a+1)<-1<0\)，故 \(f(x)\) 在 \((-\infty,a)\) 上单调递减. 当 \(x\geqslant a\) 时，\(|x-a|=x-a\)，所以
\[
f(x)=x^2-(2a-1)x
\]
其导数为 \(2x-(2a-1)\geqslant1>0\)，故 \(f(x)\) 在 \([a,+\infty)\) 上单调递增.

\item 令 \(q(x)=f(x)+\frac4x\)，只在 \(x>0\) 内讨论. 当 \(0<x<a\) 时，由上面的分段式，\(q'(x)=2x-(2a+1)-\frac4{x^2}<0\)，所以 \(q(x)\) 严格递减，且 \(q(0+)=+\infty\). 在 \(x=a\) 处，
\[
q(a)=-a(a-1)+\frac4a=-\frac{(a-2)(a^2+a+2)}{a}
\]
当 \(a=2\) 时，\(q(a)=0\)，因此 \(q(x)>0\)（\(0<x<a\)），这一段没有零点；当 \(a>2\) 时，\(q(a)<0\)，由连续性和严格递减性知 \((0,a)\) 内恰有一个零点.

当 \(x\geqslant a\) 时，\(q(x)=x^2-(2a-1)x+\frac4x\)，并且
\(q'(x)=2x-(2a-1)-\frac4{x^2}\)，\(q''(x)=2+\frac8{x^3}>0\).
又 \(q'(a)=1-\frac4{a^2}\geqslant0\)，所以 \(q\) 在 \([a,+\infty)\) 上严格递增，且 \(q(x)\to+\infty\). 当 \(a=2\) 时，\(q(a)=0\)，故这一段只有端点零点 \(x=2\)；当 \(a>2\) 时，\(q(a)<0\)，故这一段恰有一个零点. 综上，当 \(a=2\) 时有 1 个零点，即 \(x=2\)；当 \(a>2\) 时有 2 个零点.
\end{enumerate}
\end{solution}
